\documentclass[10pt,a4paper]{amsart}
\usepackage[margin=19mm]{geometry}
\usepackage{amsmath,amssymb,mathtools}
\usepackage[scr=rsfs,cal=euler]{mathalfa}
\usepackage{xfrac}
\usepackage[unicode,hidelinks,bookmarks=false]{hyperref}
\newcommand{\ie}{i.e.\ }
\newcommand{\cf}{cf.\ }
\newcommand{\resp}{respectively\ }
\newcommand{\Subsection}{\subsection}
\providecommand{\nobreakdash}{\nobreak\hbox{-}}
\setlength{\emergencystretch}{2em}
\multlinegap0pt
\usepackage{amssymb}
\newtheorem{proposition}{Proposition}
\newtheorem{lemma}{Lemma}
\newcommand{\C}{{\mathbb{C}}}
\title{Characteristic functions of two meromorphic functions weakly sharing three small\\ functions with bi-weights}
\author{Si Duc Quang, Phung Nguyen Ngoc Anh}
\date{Source: arXiv:2605.26996v1 (CC BY 4.0)}
\begin{document}
\maketitle

\noindent\emph{Source and licence.} Characteristic functions of two meromorphic functions weakly sharing three small\\ functions with bi-weights. Version arXiv:2605.26996v1 (CC BY 4.0), distributed by arXiv under the Creative Commons Attribution 4.0 International licence, http://creativecommons.org/licenses/by/4.0/, as recorded in the arXiv licence metadata for this exact version and shown on the abstract page https://arxiv.org/abs/2605.26996v1. Full licence notice: \texttt{licenses/arxiv-2605.26996-CC-BY-4.0.txt}.\par\smallskip

\noindent\textit{Editorial scope.} Counted unit: the article's lemma 3.1 (source0.tex lines 241--249). The article's complete original proof of it is delivered with it (source0.tex lines 251--390, one \texttt{proof} environment of 140 lines). No mathematics is written, repaired or re-derived: the statement and the proof are the article's own lines, in the article's own notation, and no retained line is substituted or re-flowed.  Retained uncounted as the source-local context the proof needs: lines 181--186; lines 195--237.  Nothing was omitted from any retained span, so the package carries no omission record.  The retained lines cite 3 works, whose entries are transcribed verbatim from the article's own bibliography source in the e-print and delivered with short editorial labels recorded in the item's reference mapping.  No retained line contains a figure environment, an image-inclusion command or drawing code, so no image is delivered and the package declares no asset dependency.  Editorial formatting only, in addition: an AMS \texttt{amsart} preamble that restates the article's own declarations and macros used by the retained lines, namely the environments \texttt{proposition}, \texttt{lemma} on separate counters, together with the standard packages those lines need.  The retained environments print in this document as Proposition 1, Lemma 1 in order of appearance, whereas the article prints the same items with the numbering of its own full document.  The complete native source of the article as distributed in the arXiv e-print for this exact version is bundled unchanged as \texttt{sources/201460/source0.tex}.
\begin{proposition}[\cite{QA}, Proposition 3.8]\label{2.2}
Let $f$ and $g$ be non-constant meromorphic functions and $a_i,b_i\ (i=1,2,3)$ $ (a_i\ne a_j, b_i\ne b_j, i\ne j)$ be small functions (with respect to $f$ and $g$). Assume that $f$ is not a quasi-M\"{o}bius transformation of $g$. Then for every positive integer $n$ we have the following inequality
$$\|\  N(r,\nu)\le \overline N(r,|\nu^0_{f-a_1}-\nu^0_{g-b_1}|)+\overline N(r,|\nu^0_{f-a_2}-\nu^0_{g-b_2}|)+S(r), $$
where $S(r)=o(T(r,f)+T(r,g))$ outside a finite Borel measure set of $[1,+\infty)$ and $\nu$ is the divisor defined by 
$\nu (z)=\max\{0,\min\{\nu^0_{f-a_3}(z),\nu^0_{g-b_3}(z)\}-1\}.$
\end{proposition}

\begin{lemma}\label{2.6}
Let $f$ and $g$ be non-constant meromorphic functions and $a_i\ (i=1,2,3)$ $ (a_i\ne a_j, i\ne j)$ small functions (with respect to $f$ and $g$). Let $n_1,n_2,n_3,k$ be positive integers ($k$ may be $+\infty$) and $\Delta=n_1n_2n_3-n_1-n_2-n_3-2>0$. Suppose that $f$ and $g$ weakly share three small functions $a_i\ (1\le i\le 3)$ with the bi-weight $(n_i,k)$. If $f$ is not a quasi-M\"{o}bius transformation of $g$ then
$$\|\ \Delta\overline N_{(n_u+1,k)}(r,\nu^{a_u}_f)\le \frac{2(n_v+1)(n_w+1)}{k+1}T(r)+S(r),$$
for every permutation $(u,v,w)$ of $\{1,2,3\}$, where $T(r)=T(r,f)+T(r,g)$, $S(r)=o(T(r))$.
\begin{proof}
Denote by $S$ the discrete subset of $\overline\C$ with $\|\ N(r,S)=o(T(r,f)+T(r,g))$ such that $E^S_{n_i,k)}(a_i,f)=E^S_{n_i,k)}(a_i,g)\ (i=1,2,3).$ By using a quasi-M\"{o}bius transformation if necessary, we may assume that $a_1=0,a_2=\infty,a_3=1$. Without loss of generality, we prove the lemma in the case of $u=1$. We set
\begin{align*}
H=\dfrac{f'}{f-1}-\dfrac{g'}{g-1}
\end{align*}
Since $f$ is not a M\"{o}bius transformation of $g$, $H\not\equiv 0$. Let $z\ (z\not\in S)$ be a zero of $f$ with $\nu^0_f(z)\le k$. If $\nu^0_f(z)> n_1$ then $z$ must be a zero of $H$ with multiplicity at least $n_1$. Therefore, we have
\begin{align}\label{3.6}
\begin{split}
n_1\overline N_{(n_1+1,k)}(r,\nu^{a_1}_f)&\le N(r,\nu^0_{H})\le T(r,H)=N(r,\nu^{\infty}_{H})+S(r)\\
&\leq \overline N(r,|\nu^{\infty}_f-\nu^{\infty}_g |)+\overline N(r,|\nu^{1}_f-\nu^{1}_g |)+S(r)\\
&\leq \sum_{i=2,3}(\overline N_{(n_i+1,k)}(r,\nu^{a_i}_{f})+\sum_{h=f,g}\overline N_{(k+1}(r,\nu^{a_i}_{h}))+S(r)\\
&\leq \sum_{i=2,3}\overline N_{(n_i+1,k)}(r,\nu^{a_i}_{f})+\dfrac{2}{k+1}T(r)+S(r).
\end{split}
\end{align}
Here, the third inequality comes from Proposition \ref{2.2}. Because $a_1,a_2,a_3$ play the same role, we have
$$n_2\overline N_{(n_2+1,k)}(r,\nu^{a_2}_f)\leq\sum_{i=1,3} \overline N_{(n_i+1,k)}(r,\nu^{a_i}_{f})+\dfrac{2}{k+1}T(r)+S(r).$$
Combining (\ref{3.6}) and the above inequality, we get
\begin{align*}
n_1n_2\overline N_{(n_1+1,k)}(r,\nu^{a_1}_f)\le&(n_2+1)\overline N_{(n_3+1,k)}(r,\nu^{a_3}_{f})\\
&+\overline N_{(n_1+1,k)}(r,\nu^{a_1}_{f})+\dfrac{2n_2+2}{k+1}T(r)+S(r).
\end{align*}
Thus,
\begin{align*}
(n_1n_2-1)\overline N_{(n_1+1,k)}(r,\nu^{a_1}_f)\le (n_2+1)\left(\overline N_{(n_3+1,k)}(r,\nu^{a_3}_{f})+\dfrac{2}{k+1}T(r)\right)+S(r).
\end{align*}
Similarly, we have
\begin{align*}
(n_3n_2-1)\overline N_{(n_3+1,k)}(r,\nu^{a_3}_f))\le (n_2+1)\left(\overline N_{(n_1+1,k)}(r,\nu^{a_1}_{f})+\dfrac{2}{k+1}T(r)\right)+S(r).
\end{align*}
From above two inequalities, we obtain
\begin{align*}
(n_3n_2-1)(n_1n_2-1)&\overline N_{(n_1+1,k)}(r,\nu^{a_1}_f)\le (n_3n_2-1)\dfrac{2(n_2+1)}{k+1}T(r)\\ 
&+(n_2+1)^2\left(\overline N_{(n_1+1,k)}(r,\nu^{a_1}_{f})+\dfrac{2}{k+1}T(r)\right)+S(r).
\end{align*}
This implies that
$$\Delta\overline N_{(n_1+1,k)}(r,\nu^{a_1}_f)\le \frac{2(n_2+1)(n_3+1)}{k+1}T(r)+S(r).$$
The lemma is proved.
\end{proof}
\end{lemma}

\begin{lemma}\label{3.1} 
    \textit{Let $f$ and $g$ be non-constant meromorphic functions, $f$ is not a M\"{o}bius transformation of $g$. Assume that $f$ and $g$ weakly share three values $a_1=0, a_2=1, a_3=\infty$ with the bi-weights $(n_i,k)$ respectively, where $n_1,n_2,n_3$ and $k$ be positive integers or $+\infty$ with $n_i\le k\ (1\le i\le 3)$ satisfying $\Delta:=n_1n_2n_3-n_1-n_2-n_3-2>0$. We have one of the followings holds:
\begin{itemize}
    \item [(a)] $T(r,f)\le T(r,g)+\left(\displaystyle\frac{8\sum_{1\le i<j\le 3}(n_i+1)(n_j+1)}{\Delta(k+1)}+\frac{5}{k+1}\right)T(r) +S(r)$.
    \item [(b)]$
    T(r, f)+T(r, g)\le N(r,g)+N\left(r, \dfrac{1}{g-1}\right) +N\left(r, \dfrac{1}{g}\right)+N_0(r)\\
    +\displaystyle\left(\frac{116\sum_{1\le i<j\le 3}(n_i+1)(n_j+1)}{\Delta(k+1)}+\frac{78}{k+1}\right)T(r)+ S(r).$
\end{itemize}}
\end{lemma}

\begin{proof}. 
We define $h_1:=\dfrac{f}{g}$ and $h_2:=\dfrac{f-1}{g-1}$. It easy to see that $T(r, h_i)\le T(r)$, $i=1,2$. Because $f$ is not a M\"{o}bius transformation of $g$, $\beta =\dfrac{h_1'}{h_1}-\dfrac{h_2'}{h_2} \not\equiv 0$. Set $\alpha=-\dfrac{1}{\beta}\dfrac{h_2'}{h_2}$. From Lemma \ref{2.6}, we have
\begin{align}\label{3.2}
\begin{split}
    N(r, \beta)&\le N(r, \nu_1)\le \sum_{h=f, g} \sum_{i=1}^{3}\left[\overline{N}_{(n_i+1,k)}(r, \nu_h^{a_i})+\overline{N}_{(k+1}(r, \nu_h^{a_i})\right]\\
    &\le \left(\frac{4\sum_{1\le i<j\le 3}(n_i+1)(n_j+1)}{\Delta(k+1)}+\frac{3}{k+1}\right)T(r)+S(r),\\
\end{split} 
\end{align} 
\begin{align}\label{3.3}\begin{split}
    T(r, \alpha)&\le T\left(r, \dfrac{h_2'}{h_2}\right)+T(r, \beta)=N\left(r, \dfrac{h_2'}{h_2}\right)+N(r, \beta)+S(r)\\
    & \le \sum_{h=f, g} \sum_{i=2,3}\left(\overline{N}_{(n_i+1,k)}(r, \nu_h^{a_i}) +\overline{N}_{(k+1}(r, \nu_h^{a_i})\right)+ N(r,\beta) +S(r)\\
    & \le \left(\frac{8\sum_{1\le i<j\le 3}(n_i+1)(n_j+1)}{\Delta(k+1)}+\frac{5}{k+1}\right)T(r) +S(r).
\end{split}\end{align}
where $\nu_1$ is the reduced divisor counting all points $z$ such that 
$$\nu_f(z)\ne \nu_g(z)\text{ or }\nu_{f-1}(z)\ne \nu_{g-1}(z).$$

It is easy to compute that
\begin{align*}
    f = \dfrac{h_1-h_1h_2}{h_1-h_2}=\dfrac{h_2^{-1}-1}{h_2^{-1}-h_1^{-1}}
\end{align*}
Thus 
\begin{align*}
    f-\alpha=\dfrac{(1-\alpha)h_2^-1+\alpha h_1^{-1}-1}{h_2^{-1}-h_1^{-1}}
\end{align*}
Now, we set $F_1=(\alpha-1)h_2^{-1}-\alpha h_1^{-1}+1, F_2=(1-\alpha)h_2^{-1}, F_3=\alpha h_1^{-1}$. Then
$$F_1+F_2+F_3=1$$
and
\begin{align*}
    f-\alpha=\dfrac{F_1}{h_2^{-1}-h_1^{-1}}
\end{align*}
By using Claim 3.11 in \cite{QK}, we have
\begin{align}\label{3.4}
\begin{split}
 N\left(r,\dfrac{1}{f-\alpha}\right)&-N(r,f-\alpha)\\
&\ge N\left(r,\dfrac{1}{F_1}\right)-N\left(r,\dfrac{1}{h_2^{-1}-h_1^{-1}}\right)-N(r,\alpha)\\
&\ge N\left(r,\dfrac{1}{F_1}\right)-N\left(r,\dfrac{1}{h_2^{-1}-h_1^{-1}}\right)+N(r,f)-2T(r,\alpha).
\end{split}
\end{align}
\noindent We consider the following two cases:

\textbf{Case 1:} $F_1,F_2,F_3$ are linearly dependent. Then there exist constants $c_1,c_2,c_3$, not all zeros, such that
\begin{align*}
    c_1F_1+c_2F_2+c_3F_3=0
\end{align*}
\begin{itemize}
    \item If $c_1 =0,$ then $c_2F_2+c_3F_3=0$. This implies that 
		$$\frac{f-1}{f}=-\frac{c_2}{c_3}\frac{1-\alpha}{\alpha}\frac{g-1}{g}.$$
		Therefore
		$$T(r,f)\le T(r,g)+T(r,\alpha)+O(1).$$
	\item If $c_1\neq 0$ then $c_2'h_2+c_3'h_3\equiv 1$, where $\displaystyle c_2'=(1-\frac{c_2}{c_1}),c_3'=(1-\frac{c_3}{c_1}).$
		It is easy to show that
		$$f(f-1)=c_2'(1-\alpha)(g-1)f+c_3'\alpha(f-1).$$
		Thus
		$$2T(r,f)\le T(r,g)+T(r,f)+T(r,\alpha)+O(1),$$
		\text{ i.e.,}$$\ T(r,f)\le T(r,g)+T(r,\alpha)+O(1).$$
\end{itemize}     
Combining this inequality with the inequality (\ref{3.3}), we obtain the desired inequality (a) of the lemma.

\textbf{Case 2:} $F_1,F_2,F_3$ are linearly independent. By inequality (1.1.18) in \cite[Lemma 1.102]{PPC}, we have
 \begin{align*}
			T(r,F_1)\le N\Big(r,\frac{1}{F_1}\Big)+2\sum_{i=1,2}\left(\overline{N}(r,F_i)+\overline{N}(r,\frac{1}{F_i})\right)+6T(r,\alpha)+S(r).
\end{align*}
Combining this with inequality (\ref{3.3}), thus
\begin{align*}
			N\Big(r,\dfrac{1}{f-\alpha}\Big)&\ge T(r,(h_1)-N\Big(r,\dfrac{1}{h_2^{-1}-h_1^{-1}}\Big)+N(r,f)\\
			&-2\sum_{i=1,2}\left(\overline{N}(r,h_i)+\overline{N}(r,\frac{1}{h_i})\right)-8T(r,\alpha)+S(r).
\end{align*}	

We define $F:=\dfrac{\alpha-1}{h_2}-\dfrac{\alpha}{h_1}$ and $G:=\dfrac{1}{h_2}-\dfrac{1}{h_1}$. Then we have 
\begin{align}\label{3.13}
F'=\Big(\alpha'-(\alpha-1)\dfrac{F_2'}{F_2}\Big)\dfrac{1}{F_2}-\Big(\alpha'-\alpha\dfrac{F_1'}{F_1}\Big)\dfrac{1}{F_1}.
\end{align}
Let $\gamma=-\alpha'+\alpha(\alpha-1)\Big(\dfrac{h_2'}{h_2}-\dfrac{h_1'}{h_1}\Big).$

$\bullet$  If $\gamma\equiv 0$ then there is a constant $c$ such that
		$\frac{h_2}{h_1}=\frac{c(\alpha-1)}{\alpha}$, i.e, $\frac{f-1}{f}=\frac{c(\alpha-1)}{\alpha}\frac{g-1}{g}$. It yields that
		$$T(r,f)\le T(r,g)+T(r,\alpha).$$
		Combining this inequality with inequality (\ref{3.3}), we again obtain the desired inequality (a) of the lemma.

$\bullet$ If $\gamma \not\equiv 0$, then it follows from (\ref{3.13}) and the definition of $F$ that
$$\frac{1}{F_2}=\frac{1}{\gamma}\left(\Big(\alpha'-\alpha\dfrac{F_1'}{F_1}\Big)F-\alpha F'\right)\text{ and } \frac{1}{F_1}=\frac{1}{\gamma}\left(\Big(\alpha'-(\alpha-1)\dfrac{F_2'}{F_2}\Big)F-(\alpha-1)F'\right).$$
This implies that
$$G=\frac{1}{\gamma}\left((\alpha-1)(\frac{F_2'}{F_2}-\alpha\frac{F_1'}{F_1})F+F'\right).$$
We follows the argument in \cite{QK} to give the following estimate. First, we have
\begin{align*}
T(r,G)&\le T(r,\gamma)+T\left(r,(\alpha-1)(\frac{F_2'}{F_2}-\alpha\frac{F_1'}{F_1})F+F')\right)\\
&\le N(r,\gamma)+m(r,F)+N\left(r,(\alpha-1)(\frac{F_2'}{F_2}-\alpha\frac{F_1'}{F_1})F+F')\right)+S(r)\\
&\le N(r,\gamma)+m(r,F)+N(r,F)+\overline N(r,F)+2N(r,\alpha)+N(r,\nu_1)+S(r)\\
&\le T(r,F)+2N(r,\nu_1)+4N(r,\alpha)+S(r),
\end{align*}
Hence, we obtain
\begin{align*}
N\Big(r,\dfrac{1}{f-\alpha}\Big) &\ge T(r,F_2^{-1}-F_1^{-1})-N\Big(r,\dfrac{1}{F_2^{-1}-F_1^{-1}}\Big)+N(r,f)-2N(r,\nu_1)\\
&-2\sum_{i=1,2}\left (\overline{N}(r,F_i)+\overline N(r,\frac{1}{F_i})\right)-4N(r,\alpha)-8T(r,\alpha)+S(r)\\
&=m\Big(r,\dfrac{1}{F_2^{-1}-F_1^{-1}}\Big)+N(r,f)-2\sum_{i=1,2}\left (\overline{N}(r,F_i)+\overline N(r,\frac{1}{F_i})\right)\\
&-2N(r,\nu_1)-12T(r,\alpha)+S(r).
\end{align*}
From $G=F_2^{-1}-F_1^{-1}$, we have $G'=\lambda_1 F_1^{-1}-\lambda_2 F_2^{-1}$ where $\lambda_1=\dfrac{F_1'}{F_1},\lambda_2=\dfrac{F_2'}{F_2}$. Obviously, $\lambda_1-\lambda_2 \neq 0$. Therefore,
$F_2^{-1}=\frac{\lambda_1 G}{\lambda_1-\lambda_2}+\frac{G'}{\lambda_1-\lambda_2},$
and hence
\begin{align*}
f&=\dfrac{F_2^{-1}-1}{F_2^{-1}-F_1^{-1}}=-\dfrac{1}{F_2^{-1}-F_1^{-1}}+\dfrac{F_2^{-1}}{F_2^{-1}-F_1^{-1}}\\
&=-\dfrac{1}{F_2^{-1}-F_1^{-1}}+\dfrac{\lambda_1}{\lambda_1-\lambda_2}+\dfrac{G'}{(\lambda_1-\lambda_2)G}.
\end{align*}
This yields that
\begin{align*}
m(r,f)&\le m\Big(r,\dfrac{1}{F_2^{-1}-F_1^{-1}}\Big)+m(r,\frac{1}{\lambda_1-\lambda_2})+S(r)\\
&\le m\Big(r,\dfrac{1}{F_2^{-1}-F_1^{-1}}\Big)+T(r,\lambda_1-\lambda_2)+S(r)\\
&\le m\Big(r,\dfrac{1}{F_2^{-1}-F_1^{-1}}\Big)+N(r,\nu_1)+S(r).
\end{align*}
It implies that
\begin{align*}
N\Big(r,\dfrac{1}{f-\alpha}\Big)&\ge T(r,f)-2\sum_{i=1,2}\left (\overline{N}(r,F_i)+\overline N(r,\frac{1}{F_i})\right)\\
&-3N(r,\nu_1)-12T(r,\alpha)+S(r).+S(r).
\end{align*}
Hence 
\begin{align*}
N\Big(r,\dfrac{1}{g'}\Big)+N_0(r)+N(r,\beta)&\ge T(r,f)-2\sum_{i=1,2}\left (\overline{N}(r,F_i)+\overline N(r,\frac{1}{F_i})\right)\\
&-3N(r,\nu_1)-12T(r,\alpha)+S(r).
\end{align*}
From the the second main theorem, we have
\begin{align*}
    T(r,f)+T(r,g)&\le T(r,f)+N(r,g)+N\left(r, \dfrac{1}{g-1}\right) +N\left(r, \dfrac{1}{g}\right)-N\left(r, \dfrac{1}{g'}\right)+S(r)\\
    &\le N(r,g)+N\left(r, \dfrac{1}{g-1}\right) +N\left(r, \dfrac{1}{g}\right)+N_0(r)+N(r, \beta)\\
    &+2\sum_{i=1,2}\left(\overline{N}(r, h_i)+\overline{N}\left(r, \dfrac{1}{h_i}\right)\right)+3N(r, \nu_1)+12T(r, \alpha)+S(r)
\end{align*}
From Lemma \ref{2.6}, we have
\begin{align*}
    \sum_{i=1,2}&\left(\overline{N}(r, h_i)+\overline{N}\left(r, \dfrac{1}{h_i}\right)\right)\le \sum_{i=1}^3\overline N(r,|\nu^{a_i}_f-\nu^{b_i}_g|)\\
    &\le \sum_{i=1}^3\overline N_{(n_i+1,k)}(r,\nu^{a_i}_f)+\sum_{i=1}^3\left(\overline N_{(k+1}(r,\nu^{a_i}_f)+\overline N_{(k+1}(r,\nu^{b_i}_g) \right)+S(r)\\
    &\le \left(\frac{2\sum_{1\le i<j\le 3}(n_i+1)(n_j+1)}{\Delta(k+1)}+\frac{3}{k+1}\right)T(r)+S(r).
\end{align*} 
Combining the above inequalities, we get
\begin{align*}
    T(r, f)+T(r, g)&\le N(r,g)+N\left(r, \dfrac{1}{g-1}\right) +N\left(r, \dfrac{1}{g}\right)+N_0(r)\\
    &\le N(r,g)+N\left(r, \dfrac{1}{g-1}\right) +N\left(r, \dfrac{1}{g}\right)+N_0(r)\\
    &+\left(\frac{116\sum_{1\le i<j\le 3}(n_i+1)(n_j+1)}{\Delta(k+1)}+\frac{78}{k+1}\right)T(r)+ S(r). \\
\end{align*}
We obtain the assertion (b).
\end{proof}

\begin{thebibliography}{99}
\bibitem[PPC]{PPC} P.C. Hu, P.Li, C.C.Yang, \textit{Unicity of Meromorphic Mappings}, Kluwer Acdemic Publishers (2003).
\bibitem[QA]{QA} N. V. An and S. D. Quang, \textit{Two meromorphic functions sharing four pairs of small functions}, Bull. Korean Math. Soc. \textbf{54} (2017) 1159--1171.
\bibitem[QK]{QK} S. D. Quang and N. X. Ky, \textit{On the characteristic functions of meromorphic functions sub-weighted sharing three values}, J. Math. Anal. Appl. Volume \textbf{538}, Issue 1, (2024) 128356.
\end{thebibliography}
\end{document}
